\chapter{Class Diagram}

Di seguito viene presentato il diagramma delle classi della struttura 
\texttt{Graph}, realizzato seguendo la notazione UML.

\section{Struttura Generale}

Il sistema è composto da tre elementi principali:

\begin{itemize}
    \item la struttura \texttt{Edge}, che rappresenta una tratta 
    del grafo;
    \item la struttura \texttt{State}, che rappresenta uno stato 
    nella coda di priorità della variante di Dijkstra;
    \item la classe \texttt{Graph}, che incapsula il grafo e 
    l'algoritmo di ricerca del percorso ottimo.
\end{itemize}

\section{Diagramma UML}

\begin{figure}[ht]
\centering
\includegraphics[width=0.85\textwidth]{images/classdiagram.png}
\caption{Class diagram UML di Percorsi nell'Arcipelago}
\label{fig:classdiagram}
\end{figure}

\section{Relazioni tra le Classi}

\begin{itemize}
    \item \textbf{Graph} ha una relazione di \textbf{composizione} 
    con \texttt{Edge}: gli archi sono contenuti nella lista di 
    adiacenza \texttt{adj} e il loro ciclo di vita è interamente 
    gestito da \texttt{Graph}.
    \item \textbf{Graph} ha una relazione di \textbf{dipendenza} 
    con \texttt{State}: la struttura \texttt{State} viene creata 
    e utilizzata localmente all'interno del metodo 
    \texttt{dijkstra()}, senza essere memorizzata come attributo 
    della classe.
    \item \textbf{State} espone il metodo \texttt{effectiveCost()} 
    e l'\texttt{operator>} che permettono alla 
    \texttt{priority\_queue} di ordinare gli stati per costo 
    effettivo minimo, realizzando un min-heap.
\end{itemize}